
Combining~\cref{thm:end:end_is_c} with~\cref{thm:decomposition_in_equivalence_classes}, we obtain the desired decomposition of~$\sa$.

\begin{coro}
	Let $\alpha \vDash n$. Then $\sa = \bigoplus_{E\in \mc E(\alpha)} \se$ is a decomposition into indecomposable submodules.
\end{coro}

\begin{exam}
	\label{exa:decomposable_module}
	 In general,~\cref{thm:end:end_is_c} does not hold for skew modules $\sse$.
	The two tableaux
	\begin{align*}
	T_0 =
	\begin{ytableau}
	1 \\
	*(gray) & *(gray) & 2
	\end{ytableau}
	\overset{\pi_1}{\longrightarrow}
	T_1 = 
	\begin{ytableau}
	2 \\
	*(gray) & *(gray) & 1
	\end{ytableau}
	\end{align*}
	form an equivalence class~$E$. Let $n=2$ and $\ab=\sh(T_0)$. 
	Observe that we obtain an idempotent~$\He$-endomorphism $\varphi$ by setting $\varphi(T_0) = \varphi(T_1) = T_1$.
	 Clearly, $\varphi$ is none of the trivial idempotents $0,\id\in \End_{\He}(\sse)$. Thus, $\End_{\He}(\sse)\neq \K\id$. Moreover, we obtain a decomposition
	\begin{align*}
	\sse = \varphi(\sse) \oplus (\id -\varphi) (\sse) = \spa_\K(T_1) \oplus \spa_\K(T_0-T_1)
	\end{align*} 
	into two submodules of dimension~$1$.
	
	This example also illustrates how the argumentation of this section can fail when it is applied to skew modules. Note that $D(T_1) \subseteq D(T_0)$. So, we may try to set
	 \begin{align*}
		i &= \max \big\{ k \in [n] \mid T_1^{-1}(k) \neq T^{-1}_0(k) \big\}, \\
		j &= \min \big\{k \in [n] \mid k>i \text{ and } i\att_{T_0} k \big\}.
	\end{align*}
	as before.
	But then $i=2$ so that~$j$ does not exist. 
\end{exam}
